\documentclass[11pt]{article}

\usepackage[a4paper,margin=1in]{geometry}
\usepackage{amsmath,amssymb,amsthm,mathtools}
\usepackage[hidelinks]{hyperref}
\hypersetup{
  pdfauthor={Bangyen Pham},
  pdftitle={Coefficient Mass of Polynomial Multiples at Complex Roots},
  pdfsubject={Lower bounds for the coefficient order statistics and logarithmic coefficient mass of polynomial multiples with prescribed complex roots},
  pdfkeywords={polynomial multiples, coefficient mass, complex roots, Newton polygon, argument principle},
  pdflang={en-US}
}

\newtheorem{theorem}{Theorem}[section]
\newtheorem{lemma}[theorem]{Lemma}
\newtheorem{corollary}[theorem]{Corollary}
\newtheorem{proposition}[theorem]{Proposition}
\newcommand{\R}{\mathbb R}
\newcommand{\N}{\mathbb N}

\title{Coefficient Mass of Polynomial Multiples at Complex Roots}
\author{\shortstack{Bangyen Pham\\
\small Independent Researcher\\
\small\href{mailto:bangyenp@gmail.com}{\texttt{bangyenp@gmail.com}}\\
\small\href{https://orcid.org/0009-0006-9928-2088}{ORCID: 0009-0006-9928-2088}}}
\date{September 2026}

\begin{document}
\maketitle

\begin{abstract}
For a nonzero polynomial $F=\sum_jf_jx^j$ of degree $D$, let $b_k(F)$ be the
$k$-th largest nonleading coefficient magnitude and
$\Lambda(F)=\sum_j\log^+|f_j|$ its logarithmic coefficient mass.  For real
roots $2\le r_1\le\cdots\le r_L$ a companion paper proves
$b_k(F)\ge|f_D|\prod_{i\ge k}(r_i-1)$ for every multiple of
$\prod_i(x-r_i)$, and hence a mass bound quadratic in $L$.  We ask what
survives for complex roots.  Roots in full orbits of $x\mapsto\zeta x$,
$\zeta^m=1$, in particular purely imaginary pairs $\pm ic$, reduce exactly to
real roots: the infimum of every $b_k$ and of $\Lambda$ over monic multiples
is the same as at the real roots $c^2$.  Without such structure only the
first row survives: $b_1(F)\ge|f_D|\prod_j(|\alpha_j|-1)$ at arbitrary roots
of modulus greater than one, while $b_2$ and every superlinear mass bound
fail for pairs within any fixed angle of the imaginary axis.  Rows and mass
return, up to constants, when the moduli are separated by a factor $9$, by
an Archimedean Newton polygon argument.  Real roots of both signs charge one
root of each sign per excluded coefficient, with the sharp constant
$\frac14$ in the mass.  For $K$ roots of modulus at least $R$ in a sector of
half-width $\delta$, the mass of multiples of degree $O(K)$ is of order
$\min(K^2,K/\delta)\log R$, with both constants attained in the limit; at
Gaussian-integer roots the degree restriction is removed for integer
multiples.
\end{abstract}

\noindent\textbf{Keywords.} Polynomial multiples; coefficient mass; complex
roots; Newton polygon; argument principle.

\noindent\textbf{2020 Mathematics Subject Classification.} Primary 11C08; Secondary 30C15, 26C10.

\section{Introduction}

For a nonzero $F(x)=\sum_jf_jx^j$ of degree $D$ write
\[
  \Lambda(F)=\sum_{j:f_j\ne0}\log^+|f_j|,\qquad \log^+x=\max(0,\log x),
\]
for its logarithmic coefficient mass, and $b_k(F)$ for the $k$-th largest of
the nonleading magnitudes $|f_j|$, $j<D$.  For real roots
$2\le r_1\le\cdots\le r_L$, repetitions allowed, the companion
paper~\cite{coefficient-mass} proves that every nonzero real or complex
multiple $F$ of $\prod_{i=1}^L(x-r_i)$ satisfies
\begin{equation}\label{eq:order}
  b_k(F)\ \ge\ |f_D|\prod_{i=k}^L(r_i-1)\quad(1\le k\le L),
  \qquad
  \Lambda(F)\ \ge\ (L+1)\log|f_D|+\sum_{i=1}^Li\log(r_i-1)
\end{equation}
\cite[Theorem~3.2, Corollaries~3.3 and~3.4]{coefficient-mass}, uniformly in
the degree; we call the first bound at a given $k$ the row $k$.  The row
$k=1$ needs only $r_1>1$.  This paper asks what remains of
\eqref{eq:order} when the prescribed roots are complex.

Reduced height and reduced length, the least height and length of a
normalized multiple, have been studied for real and complex
coefficients~\cite{dj,schinzel-complex}, and power series with restricted
coefficients and a root on a given ray in~\cite{bbbp-ray}; the first row
below is the complex form of the bound $H(P)\ge|P(1)|$ of~\cite{dj}.

\paragraph{Exact reductions.}  If the prescribed factor has the form
$R(x^m)$, multisection reduces every multiple to a multiple of $R$ with a
subset of the coefficients (Theorem~\ref{thm:transfer}).  In particular a
purely imaginary pair $\pm ic$ is worth exactly one real root $c^2$, and all
the real-root results transfer (Corollary~\ref{cor:imag}); the same holds for
roots on a common ray (Corollary~\ref{cor:ray}).

\paragraph{Arbitrary angles.}  The first row survives at every root set
(Proposition~\ref{prop:rowone}), but nothing more does in general: pairs
within any fixed angle of the imaginary axis admit multiples with $b_2=0$
and mass linear in the number of pairs (Proposition~\ref{prop:noangle}).
Separated moduli restore the rows and the mass up to constants, through an
Archimedean Newton polygon (Theorem~\ref{thm:archnewton},
Corollary~\ref{cor:separated}).

\paragraph{Both signs.}  For real roots of both signs, Rolle's theorem on
each half-line charges one root of each sign per excluded coefficient
(Theorem~\ref{thm:bothsigns}); the resulting mass bound has the sharp
leading constant $\frac14$ (Corollary~\ref{cor:bothsignsmass},
Proposition~\ref{prop:bothsignssharp}).

\paragraph{Sectors.}  For $K$ roots of modulus at least $R$ in a sector, the
argument principle along a ray next to the sector produces many real zeros
of a real polynomial dominated coefficientwise by the multiple, to which
\eqref{eq:order} applies.  In a thin sector this gives mass of order
$K^2\log R$ for multiples of small degree (Theorem~\ref{thm:thinsector}), and
at Gaussian-integer roots, where a carry argument charges degree against
mass (Theorem~\ref{thm:gausscharge}), for all integer multiples
(Theorem~\ref{thm:thingauss}).  In a sector of half-width $\delta$ the
crossings are counted against the equidistributed share $\delta/\pi$ of the
large zeros, and the mass of multiples of degree $O(K)$ is of order
$\min(K^2,K/\delta)\log R$ (Theorem~\ref{thm:widesector},
Corollary~\ref{cor:wideconst}, Proposition~\ref{prop:widesharp}).  Whether
$\Lambda(F)\ge cK^2\log R$ holds for all integer multiples at Gaussian roots
in a wide sector is open.

\paragraph{Conventions.}
All logarithms are natural.  For $F\in\mathbb C[x]$ of degree $D$ we write
$F^*(t)=t^DF(1/t)=\sum_{d=0}^Df_{D-d}t^d$ for the reciprocal polynomial.
References of the form \cite[Theorem~3.2]{coefficient-mass} are to the
companion paper.

\section{Exact reductions}\label{sec:exact}

\begin{lemma}[Multisection]\label{lem:multisection}
Let $m\ge1$, let $R\in\mathbb C[y]$ with $R(0)\ne0$, and put $P(x)=R(x^m)$.  Write
a polynomial $F\in\mathbb C[x]$ as
\[
  F(x)=\sum_{\varepsilon=0}^{m-1}x^\varepsilon E_\varepsilon(x^m),
  \qquad E_\varepsilon(y)=\sum_{i\ge0}f_{mi+\varepsilon}\,y^i .
\]
If $P$ divides $F$, then $R$ divides every $E_\varepsilon$.  If $F$ and $R$
are real, so are the $E_\varepsilon$.
\end{lemma}

\begin{proof}
Let $\zeta=e^{2\pi i/m}$.  Since $P(\zeta^kx)=R(\zeta^{km}x^m)=P(x)$, the
polynomial $P$ divides $F(\zeta^kx)$ for every $k$, hence divides
\[
  \frac1m\sum_{k=0}^{m-1}\zeta^{-k\varepsilon}F(\zeta^kx)
  =x^\varepsilon E_\varepsilon(x^m).
\]
As $R(0)\ne0$, $P(0)\ne0$, so $x^\varepsilon$ is coprime to $P$ and
$R(x^m)\mid E_\varepsilon(x^m)$.  Divide in $\mathbb C[y]$:
$E_\varepsilon=RQ+S$ with $\deg S<\deg R$.  Then $R(x^m)$ divides
$S(x^m)$, whose degree $m\deg S$ is smaller than $m\deg R$, so $S=0$.
Reality of $E_\varepsilon$ is clear from its coefficients.
\end{proof}

\begin{theorem}[Multisection transfer]\label{thm:transfer}
Let $m$, $R$ and $P=R(x^m)$ be as in Lemma~\ref{lem:multisection}, and let
$F$ be a nonzero multiple of $P$ of degree $D$.  For the residue
$\varepsilon\equiv D\pmod m$, the polynomial $E_\varepsilon$ is a multiple of
$R$ of degree $(D-\varepsilon)/m$ with leading coefficient $f_D$, and
\[
  b_k(F)\ge b_k(E_\varepsilon)\quad(k\ge1),\qquad
  \Lambda(F)\ge\Lambda(E_\varepsilon).
\]
Conversely, for every multiple $G$ of $R$ the polynomial $G(x^m)$ is a
multiple of $P$ with the same nonzero coefficients.  Consequently, for
$\Phi=b_k$ (any fixed $k$) and for $\Phi=\Lambda$,
\[
  \inf\{\Phi(F):F\text{ monic},\ P\mid F\}
  =\inf\{\Phi(G):G\text{ monic},\ R\mid G\},
\]
over complex multiples, and also over real multiples when $R$ is real.
\end{theorem}

\begin{proof}
By Lemma~\ref{lem:multisection}, $R\mid E_\varepsilon$.  The top coefficient
of $F$ sits at $D=m\cdot\frac{D-\varepsilon}m+\varepsilon$, so it is the
coefficient of $y^{(D-\varepsilon)/m}$ in $E_\varepsilon$, and no higher
coefficient of $E_\varepsilon$ is nonzero.  The nonleading coefficients of
$E_\varepsilon$ are the $f_j$ with $j\equiv\varepsilon$, $j<D$: a
sub-multiset of the nonleading coefficients of $F$.  The $k$-th largest
element of a sub-multiset is at most the $k$-th largest of the whole, and
$\Lambda$ is a sum of nonnegative terms over coefficients; this gives the
two inequalities.  Monicity passes from $F$ to $E_\varepsilon$, which gives
$\ge$ between the infima.  For the converse, $R\mid G$ gives
$R(x^m)\mid G(x^m)$, and $G(x^m)$ has the coefficients of $G$ spread out
with zeros between them, so $b_k$ and $\Lambda$ are unchanged, and
$G(x^m)$ is monic when $G$ is.  This gives $\le$.
\end{proof}

The case $m=2$, $R(y)=\prod_j(y+c_j^2)$ is the case of purely imaginary
conjugate pairs; the sign change $y\mapsto-y$ turns $R$ into a product of
linear factors at the positive roots $c_j^2$ without changing any
coefficient magnitude.

\begin{corollary}[Purely imaginary pairs]\label{cor:imag}
Let $c_1\le\cdots\le c_K$ be positive reals, repetitions allowed, put
$P_{\rm im}=\prod_{j=1}^K(x^2+c_j^2)$ and
$P_{\rm re}=\prod_{j=1}^K(y-c_j^2)$, and let $F\in\mathbb C[x]$ be a nonzero
multiple of $P_{\rm im}$.
\begin{enumerate}
\item If $c_1^2\ge2$, then for $1\le k\le K$,
\[
  b_k(F)\ \ge\ |f_D|\prod_{j=k}^K(c_j^2-1),
\]
and the $k$ magnitudes may be taken among the positions $j\equiv D\pmod2$.
The row $k=1$ holds for $c_1>1$.
\item If $c_1^2\ge2$ and $|f_D|\ge1$, then
\[
  \Lambda(F)\ \ge\ (K+1)\log|f_D|+\sum_{j=1}^Kj\log(c_j^2-1).
\]
\item For $\Phi=b_k$ and $\Phi=\Lambda$, the infimum of $\Phi$ over monic
(real or complex) multiples of $P_{\rm im}$ equals its infimum over monic
multiples of $P_{\rm re}$.  In particular, with $c_1^2\ge2$ and $K\ge2$,
\[
  \inf_{Q\text{ monic}}\Lambda(P_{\rm im}Q)\ \asymp\
  K^2+\sum_{j=1}^Kj\log(c_j^2-1)
\]
with the absolute constants of \cite[Corollary~4.3]{coefficient-mass};
$\Lambda(F)\ge K^2/80$ for monic $F$ when $K\ge8$; and the row $k=2$ is an
infimum at $(c_1^2,c_2^2)=(2,3)$.
\end{enumerate}
\end{corollary}

\begin{proof}
Apply Theorem~\ref{thm:transfer} with $m=2$ and $R(y)=\prod_j(y+c_j^2)$;
$R(0)\ne0$.  The multiple $E_\varepsilon$ of $R$ has leading coefficient
$f_D$, and $\tilde E(y)=E_\varepsilon(-y)$ is a multiple of $P_{\rm re}$
with leading coefficient $\pm f_D$ and the same coefficient magnitudes.  \cite[Theorem~3.2]{coefficient-mass}, with \cite[Corollary~3.3]{coefficient-mass} for complex $F$, applied to
$\tilde E$ at the roots $c_j^2\ge2$ gives item~1, since
$b_k(F)\ge b_k(\tilde E)$ and the positions of $\tilde E$ are those
$j\equiv D\pmod 2$ of $F$; the clause $c_1>1$ is the $u=0$ clause of
\cite[Theorem~3.2]{coefficient-mass}.  Item~2 is \cite[Corollary~3.4]{coefficient-mass} applied to $\tilde E$, together with
$\Lambda(F)\ge\Lambda(\tilde E)$.  Item~3 is the equality of infima in
Theorem~\ref{thm:transfer}, composed with the magnitude-preserving
bijection $G(y)\mapsto\pm G(-y)$ between monic multiples of $R$ and of
$P_{\rm re}$; the three quoted statements are \cite[Corollary~4.3]{coefficient-mass},
\cite[Proposition~4.2]{coefficient-mass} and \cite[Proposition~4.4]{coefficient-mass} at
the roots $c_j^2$.
\end{proof}

Thus a purely imaginary pair $\pm ic$ is worth exactly one real root $c^2$:
the pair has two roots of modulus $c$, and the mass charges $\log c^2$ once
at its rank.  This is the right count.  By item~3 the product
$P_{\rm im}$ itself is within $O(K^2)$ of the infimum
(\cite[Corollary~3.7]{coefficient-mass}), and $\Lambda(P_{\rm im})=\Lambda(P_{\rm re})$; so a
bound of the form $\sum_{i\le2K}i\log\rho_{(i)}$ over the $2K$ roots
counted separately, which would be about twice as large, is false already
for $P_{\rm im}$.

The same proof applies to every factor of the form $R(x^m)$ with $R$ real
and all roots of $R$ of one sign: the real roots $\pm r$ in pairs
($R=\prod(y-r_j^2)$), or the $m$-th roots of $\pm s_j$ for fixed $m$.  One
more reduction needs no multisection.

\begin{corollary}[Roots on a common ray]\label{cor:ray}
Let $\theta\in\R$ and $\rho_1\le\cdots\le\rho_K$ with $\rho_1\ge2$, and let
$F\in\mathbb C[x]$ be a nonzero multiple of $\prod_j(x-\rho_je^{i\theta})$.  Then
$b_k(F)\ge|f_D|\prod_{j\ge k}(\rho_j-1)$ for $1\le k\le K$.  In
particular this holds for a real $F$ divisible by
$\prod_j(x^2-2\rho_j\cos\theta\,x+\rho_j^2)$.
\end{corollary}

\begin{proof}
$G(x)=F(e^{i\theta}x)$ has coefficients $f_je^{ij\theta}$, of the same
magnitudes, and is divisible by $\prod_j(x-\rho_j)$; apply
\cite[Corollary~3.3]{coefficient-mass}.
\end{proof}

A common ray charges each pair once, with $\log(\rho-1)$ in place of
$\log(\rho^2-1)$; the imaginary axis is special because it carries both
members of each pair.


\section{Arbitrary angles}\label{sec:angles}

\begin{proposition}[The first row at arbitrary roots]\label{prop:rowone}
Let $\alpha_1,\ldots,\alpha_N\in\mathbb C$, repetitions allowed, with
$\rho_j=|\alpha_j|$, and let $F\in\mathbb C[x]$ be a nonzero multiple of
$\prod_j(x-\alpha_j)$.
\begin{enumerate}
\item If every $\rho_j>1$, then $b_1(F)\ge|f_D|\prod_{j=1}^N(\rho_j-1)$.
\item For every $0<r<1$,
$b_1(F)\ge|f_D|\,\frac{1-r}{r}\bigl(r^N\mathcal M-1\bigr)$, where
$\mathcal M=\prod_j\max(1,\rho_j)$; with $r=N/(N+1)$,
$b_1(F)\ge|f_D|\,(\mathcal M/e-1)/N$.
\end{enumerate}
\end{proposition}

\begin{proof}
Let $F^*(t)=t^DF(1/t)=\sum_{d=0}^Df_{D-d}t^d$ and $\beta_j=1/\alpha_j$
where $\alpha_j\ne0$.

\emph{Item 1.}  Here every $\alpha_j\ne0$, and if $(x-\alpha)^\mu\mid F$
then $F^*(t)=(1-\alpha t)^\mu\,t^{D-\mu}H(1/t)$ with $F=(x-\alpha)^\mu H$, so
$\chi(t)=\prod_j(t-\beta_j)$ divides $F^*$.  Put
$B(s)=\prod_j(1-\beta_js)=s^N\chi(1/s)$, let $c=-[s^N]B$, and let
$v_d=[s^d]V$ for $V=(B+cs^N)/B=1+cs^N/B$.  The numerator $A=B+cs^N$ has
degree at most $N-1$, so $BV=A$ gives, for every $n\ge N$,
$\sum_{k=0}^N[s^k]B\;v_{n-k}=0$, that is,
$\sum_{i=0}^N\chi_iv_{m'+i}=0$ for every $m'\ge0$, where
$\chi=\sum_i\chi_it^i$.  Writing $F^*=\chi h$,
\[
  \sum_{d}v_df_{D-d}
  =\sum_{m'}h_{m'}\sum_i\chi_iv_{m'+i}=0 .
\]
Since $v_0=1$, $|f_D|\le b_1(F)\sum_{d\ge1}|v_d|$.  Now
$|c|=\prod_j|\beta_j|=\prod_j\rho_j^{-1}$, and the coefficients of $1/B$ are
dominated in modulus by those of $\prod_j(1-|\beta_j|s)^{-1}$, whose sum is
$\prod_j(1-\rho_j^{-1})^{-1}$.  Hence
$\sum_{d\ge1}|v_d|\le\prod_j\rho_j^{-1}(1-\rho_j^{-1})^{-1}
=\prod_j(\rho_j-1)^{-1}$.

\emph{Item 2.}  $F^*(0)=f_D\ne0$, and $F^*$ vanishes at every $\beta_j$
(counted with multiplicity, as in item 1).  Jensen's formula on $|t|\le r$
gives
\[
  \log\max_{|t|=r}|F^*|\ \ge\ \log|f_D|+\sum_{j:\rho_j>1/r}\log(r\rho_j)
  \ \ge\ \log|f_D|+\sum_{j}\log\bigl(r\max(1,\rho_j)\bigr),
\]
since each omitted or added term is $\le0$ on the right.  On the other side
$\max_{|t|=r}|F^*|\le|f_D|+b_1(F)\sum_{d\ge1}r^d
=|f_D|+b_1(F)\,r/(1-r)$.  Rearranging gives the first bound; for
$r=N/(N+1)$, $(1-r)/r=1/N$ and $r^N=(1+1/N)^{-N}\ge e^{-1}$.
\end{proof}

For real roots $r_j>1$ item 1 is the row $k=1$ of \cite[Theorem~3.2]{coefficient-mass}.  Both items
are of the size of the Mahler measure: they charge each root once, as
$\log\rho_j$, and do not give a second large coefficient.  At arbitrary
angles nothing more is true.

\begin{proposition}[No angle-uniform rows beyond the first]
\label{prop:noangle}
\leavevmode
\begin{enumerate}
\item Let $M\ge3$ be odd, $\rho>1$, and $0<\delta<\pi/2$.  The monic
polynomial $x^{2M}+\rho^{2M}$ is divisible by the product of the
$K=2\lfloor\delta M/\pi\rfloor+1$ conjugate pairs of modulus $\rho$ with
arguments $\pm(\frac\pi2+t\frac\pi M)$, $|t|\le\delta M/\pi$, all within
$\delta$ of the imaginary axis; yet
\[
  b_2(x^{2M}+\rho^{2M})=0,\qquad
  \Lambda(x^{2M}+\rho^{2M})=2M\log\rho\le\frac\pi\delta(K+1)\log\rho .
\]
\item For $m\ge2$ and $B\ge1$, the polynomial $x^{2m}+x+B$ has no real
root, so its roots form $m$ conjugate pairs, and $b_2=1$,
$\Lambda=\log B$.
\end{enumerate}
In particular, no bound $b_2(F)\ge\phi(\rho_1,\ldots,\rho_K)>0$ in terms of
the moduli of the prescribed pairs alone, not even for pairs within a fixed
angle $\delta$ of the imaginary axis, and no bound
$\Lambda(F)\ge c\,K^2\log\rho$ for $K$ pairs of modulus $\rho$ within a
fixed angle $\delta$, can hold for all monic multiples.
\end{proposition}

\begin{proof}
1.  The roots of $x^{2M}+\rho^{2M}$ are $\rho e^{i\pi(2l+1)/(2M)}$; for
$M$ odd, $l=(M-1)/2+t$ gives the argument $\frac\pi2+t\frac\pi M$, and the
conjugate root has the negated argument.  The integers $t$ with
$|t|\le\delta M/\pi$ number $K$.  The only nonleading nonzero coefficient
is $\rho^{2M}$, so $b_2=0$ and $\Lambda=2M\log\rho$.  Finally
$K\ge2\delta M/\pi-1$ gives $2M\le\pi(K+1)/\delta$.

2.  For real $x$ with $|x|\ge1$, $x^{2m}+x\ge|x|^{2m}-|x|\ge0$, and for
$|x|<1$, $x^{2m}+x>-1$; so $x^{2m}+x+B>0$ when $B\ge1$.  The coefficients are
$1,1,B$.

For the final sentence, item 1 with $\delta M\ge\pi$, so $K\ge3$, has $b_2=0$ while every
$\rho_j=\rho$ is fixed, and letting $M\to\infty$ at fixed $\delta$ gives
$K\to\infty$ with $\Lambda=O_\delta(K\log\rho)$.
\end{proof}

What the imaginary axis provides in Corollary~\ref{cor:imag} is therefore
not closeness of the angles but exact membership, which makes the
prescribed factor a polynomial in $x^2$.  Separated moduli, on the other
hand, do help, through the Archimedean Newton polygon.  For a nonzero
$F=\sum_if_ix^i\in\mathbb C[x]$ of degree $D$ and order $e$ at $0$, let
$\varphi$ be the least concave function on $[e,D]$ with
$\varphi(i)\ge\log|f_i|$ whenever $f_i\ne0$.  Its vertices
$D=v_0>v_1>\cdots>v_L=e$ are among these points, so $\varphi(v_t)=\log|f_{v_t}|$.
The edge joining $v_s$ to $v_{s-1}$ has length $\ell_s=v_{s-1}-v_s\ge1$ and
slope $-\sigma_s$; by concavity $\sigma_1>\sigma_2>\cdots>\sigma_L$, and
\[
  \log|f_{v_t}|=\log|f_D|+\sum_{s\le t}\ell_s\sigma_s\qquad(0\le t\le L).
\]
We call $\sigma_1,\ldots,\sigma_L$ the \emph{tropical roots} of $F$.  For
$u\in\R$ the vertex $v_s$ with $\sigma_{s+1}<u<\sigma_s$ (where
$\sigma_0=+\infty$ and $\sigma_{L+1}=-\infty$) is the one maximizing
$\log|f_i|+iu$; we call it \emph{active} at $u$.

\begin{theorem}[Archimedean Newton polygon]\label{thm:archnewton}
Let $F\in\mathbb C[x]$ be nonzero, with tropical roots $\sigma_1>\cdots>\sigma_L$.
\begin{enumerate}
\item Let $u\in\R$ satisfy $|u-\sigma_s|>\log3$ for every $s$, let $v$ be the
vertex active at $u$, and $r=e^u$.  Then
$|f_v|r^v>\sum_{i\ne v}|f_i|r^i$.  Hence $F$ has no root on $|z|=r$ and
exactly $v$ roots, counted with multiplicity, in $|z|<r$.
\item Every nonzero root $\beta$ of $F$ satisfies
$\bigl|\log|\beta|-\sigma_s\bigr|\le\log3$ for some $s$.
\end{enumerate}
\end{theorem}

\begin{proof}
1.  Let $v=v_s$, so $\sigma_{s+1}<u<\sigma_s$, and put $d_+=\sigma_s-u$ and
$d_-=u-\sigma_{s+1}$, both larger than $\log3$ (and infinite when $s=0$,
respectively $s=L$).  A concave function lies below the line through any of
its edges, so for $i>v$ the edge of slope $-\sigma_s$ gives
$\log|f_i|\le\varphi(i)\le\varphi(v)-\sigma_s(i-v)$, that is,
$|f_i|r^i\le|f_v|r^ve^{-(i-v)d_+}$; likewise the edge of slope $-\sigma_{s+1}$
gives $|f_i|r^i\le|f_v|r^ve^{-(v-i)d_-}$ for $i<v$.  Summing over the
distinct positive integers $|i-v|$,
\[
  \sum_{i\ne v}|f_i|r^i\ \le\ |f_v|r^v\Bigl(\frac1{e^{d_+}-1}+\frac1{e^{d_-}-1}\Bigr)
  \ <\ |f_v|r^v\Bigl(\frac12+\frac12\Bigr).
\]
So $|F(z)|\ge|f_v|r^v-\sum_{i\ne v}|f_i|r^i>0$ on $|z|=r$, and by Rouch\'e's
theorem $F$ and $f_vz^v$ have the same number of zeros in $|z|<r$, namely $v$.

2.  If $|\log|\beta|-\sigma_s|>\log3$ for every $s$, item~1 at
$u=\log|\beta|$ shows $F(\beta)\ne0$.
\end{proof}

Item~2 is the Archimedean form of the Newton polygon theorem: up to the
factor $3$, the moduli of the roots can only sit at the tropical roots.
Unlike the $q$-adic case, the roots of modulus near $e^{\sigma_s}$ need not
number $\ell_s$; $(x-1)^D$ has all its roots on $|z|=1$ but
tropical roots $\log\frac{i+1}{D-i}$ spread over $[-\log D,\log D]$.
Separated roots still force separated vertices, which is all the mass needs.

\begin{corollary}[Separated moduli]\label{cor:separated}
Let $F\in\mathbb C[x]$ be nonzero of degree $D$, and let
$\rho_1\le\cdots\le\rho_K$ be moduli of nonzero roots of $F$.
\begin{enumerate}
\item If $\rho_1\ge3$ and $\rho_{j+1}>9\rho_j$ for $1\le j<K$, there are
positions $y_1<\cdots<y_K<D$ with
\[
  |f_{y_j}|\ \ge\ |f_D|\prod_{i=j}^K\frac{\rho_i}3 .
\]
Hence $b_k(F)\ge|f_D|\prod_{i=k}^K(\rho_i/3)$ for $1\le k\le K$, and, if
$|f_D|\ge1$,
\[
  \Lambda(F)\ \ge\ (K+1)\log|f_D|+\sum_{j=1}^Kj\log\frac{\rho_j}3 .
\]
\item If moreover each $\rho_j$ is the modulus of two roots of $F$, counted
with multiplicity (for instance $\beta_j$ and $\bar\beta_j$ with $F$ real and
$\beta_j\notin\R$), and $\rho_1>27$ and $\rho_{j+1}>729\rho_j$, then
$b_k(F)\ge|f_D|\prod_{i=k}^K(\rho_i^2/729)$ for $1\le k\le K$, and, if
$|f_D|\ge1$,
\[
  \Lambda(F)\ \ge\ (K+1)\log|f_D|+\sum_{j=1}^Kj\log\frac{\rho_j^2}{729} .
\]
\end{enumerate}
\end{corollary}

\begin{proof}
1.  By Theorem~\ref{thm:archnewton}, item~2, choose for each $j$ an index
$s(j)$ with $|\sigma_{s(j)}-\log\rho_j|\le\log3$.  Since
$\log\rho_{j+1}-\log\rho_j>2\log3$,
\[
  \sigma_{s(j+1)}\ge\log\rho_{j+1}-\log3>\log\rho_j+\log3\ge\sigma_{s(j)},
\]
so $s(1)>s(2)>\cdots>s(K)\ge1$, and $y_j=v_{s(j)}$ are increasing positions
below $D$.  In
$\log|f_{y_j}|=\log|f_D|+\sum_{s\le s(j)}\ell_s\sigma_s$ every term has
$\ell_s\ge1$ and $\sigma_s\ge\sigma_{s(j)}\ge\log(\rho_j/3)\ge0$; keeping only
the terms $s=s(i)$, $i\ge j$, gives
$\log|f_{y_j}|\ge\log|f_D|+\sum_{i\ge j}\log(\rho_i/3)$.  Every factor
$\rho_i/3$ is at least $1$, so the $k$ positions $y_1,\ldots,y_k$ all carry
magnitudes at least $|f_D|\prod_{i\ge k}(\rho_i/3)$, which bounds $b_k(F)$.
For the mass, $\Lambda(F)\ge\log^+|f_D|+\sum_j\log^+|f_{y_j}|$, and
$\sum_j\sum_{i\ge j}\log(\rho_i/3)=\sum_ii\log(\rho_i/3)$.

2.  Fix $c$ with $3\log3<c<\log\rho_1$ and $2c<\log(\rho_{j+1}/\rho_j)$ for
every $j$, and let $W_j=[\log\rho_j-c,\log\rho_j+c]$.  These windows are
disjoint, lie in $(0,\infty)$, and $W_{j+1}$ lies above $W_j$.  We claim that
the tropical roots in $W_j$ have total multiplicity
$\sum_{\sigma_s\in W_j}\ell_s\ge2$.  By Theorem~\ref{thm:archnewton}, item~2,
$W_j$ contains a tropical root $\sigma_s$ with $|\sigma_s-\log\rho_j|\le\log3$.
If it contains another one, the claim holds.  Otherwise take
$0<\varepsilon<c-3\log3$ and $u_\pm=\sigma_s\pm(\log3+\varepsilon)$.  Both are
farther than $\log3$ from $\sigma_s$, and every other tropical root lies
outside $W_j$, at distance more than $c-(2\log3+\varepsilon)>\log3$ from
$u_\pm$, because $|u_\pm-\log\rho_j|\le2\log3+\varepsilon$.  No tropical root
other than $\sigma_s$ lies in $[u_-,u_+]\subset W_j$, so the vertices active
at $u_+$ and $u_-$ are $v_{s-1}$ and $v_s$, and by
Theorem~\ref{thm:archnewton}, item~1, $F$ has exactly
$v_{s-1}-v_s=\ell_s$ roots in $e^{u_-}<|z|<e^{u_+}$.  The two roots of modulus
$\rho_j$ are among them, as $|\log\rho_j-\sigma_s|\le\log3$; so $\ell_s\ge2$.

Now let $\sigma_{s(j)}$ be the smallest tropical root in $W_j$ and
$y_j=v_{s(j)}$.  In $\log|f_{y_j}|=\log|f_D|+\sum_{s\le s(j)}\ell_s\sigma_s$
every $\sigma_s$ is at least $\sigma_{s(j)}>0$, and the sum contains every
tropical root in $W_i$ for $i\ge j$; each of these is at least
$\log\rho_i-c$.  By the claim,
$\log|f_{y_j}|\ge\log|f_D|+2\sum_{i\ge j}(\log\rho_i-c)$.  As in item~1,
$b_k(F)\ge|f_D|\prod_{i\ge k}(\rho_i^2e^{-2c})$ and
$\Lambda(F)\ge(K+1)\log|f_D|+\sum_jj\log(\rho_j^2e^{-2c})$.  These hold for
every $c$ in a nonempty interval with infimum $3\log3$; letting
$c\downarrow3\log3$ gives the stated bounds, as $e^{6\log3}=729$.

\end{proof}

So when the prescribed moduli grow by a factor $9$ at each step, every
complex multiple pays $\sum_jj\log(\rho_j/3)$, one real root $\rho_j/3$ per
root, at every angle; with a factor $729$ a conjugate pair pays
$\log(\rho_j^2/729)$, which is the charge $\log(c_j^2-1)$ of
Corollary~\ref{cor:imag} up to the constant.  Proposition~\ref{prop:noangle}
shows that some separation is necessary; the constants $3$, $9$ and $729$ are
not optimized.

\section{Real roots of both signs}\label{sec:bothsigns}

\cite[Theorem~3.2]{coefficient-mass} needs its roots on one side of $0$.  Applied to $F(x)$
and to $F(-x)$ it charges a multiple with $p$ positive and $n$ negative
prescribed roots only for the larger of the two classes.  The mixed problem is
genuinely weaker: by the multisection transfer (the remark after
Corollary~\ref{cor:imag}, with $R=\prod_j(y-s_j^2)$), the $L=2m$ roots $\pm s_j$
of $\prod_j(x^2-s_j^2)$ are worth only the $m$ positive roots $s_j^2$, so at a
common modulus $s$ they cost about $\frac14L^2\log s$ instead of the
$\frac12L^2\log s$ of $L$ roots of one sign.  Rolle's theorem, used on the two
half-lines separately as in \cite[Proposition~4.2]{coefficient-mass}, gives the
matching lower bound: every excluded coefficient costs one root of each sign.

\begin{lemma}[Rolle on both half-lines]\label{lem:rolleside}
Let $h\in\R[t]$ be nonzero, let $z\ge1$ be an integer, and put
$T_zh=h-(t/z)h'$, so that $[t^d]T_zh=(1-d/z)[t^d]h$.  If $h$ has the zeros
$0<\xi_1\le\cdots\le\xi_m$, listed with multiplicity, $m\ge1$, then $T_zh$ has
zeros $\eta_1\le\cdots\le\eta_{m-1}$ in $(0,\infty)$, listed with multiplicity,
with $\xi_j\le\eta_j\le\xi_{j+1}$.  The same holds in $(-\infty,0)$ for the
absolute values: zeros $|\xi_1|\le\cdots\le|\xi_m|$ of $h$ there give zeros
$|\eta_1|\le\cdots\le|\eta_{m-1}|$ of $T_zh$ there with
$|\xi_j|\le|\eta_j|\le|\xi_{j+1}|$.
\end{lemma}

\begin{proof}
For $t\ne0$, $T_zh=-(t^{z+1}/z)(t^{-z}h)'$.  On $I=(0,\infty)$ the function
$\varphi=t^{-z}h$ is real-analytic, not identically zero, and has the zeros of
$h$ with the same multiplicities.  Let $x_1<\cdots<x_q$ be the distinct zeros,
of multiplicities $\mu_i$.  Writing $\varphi=(t-x_i)^{\mu_i}\psi$ with
$\psi(x_i)\ne0$ shows that $\varphi'$ vanishes to order $\mu_i-1$ at $x_i$,
and Rolle's theorem gives a zero $y_i\in(x_i,x_{i+1})$ for $i<q$.  List
$\mu_1-1$ copies of $x_1$, then $y_1$, then $\mu_2-1$ copies of $x_2$, then
$y_2$, and so on: these are $m-1$ zeros of $\varphi'$ counted with
multiplicity, the $j$-th lies in $[\xi_j,\xi_{j+1}]$, and since
$-t^{z+1}/z$ has no zero in $I$ they are zeros of $T_zh$.  For
$(-\infty,0)$ apply this to $h(-t)$, using $(T_zh)(-t)=T_z\bigl(h(-t)\bigr)$.
\end{proof}

\begin{theorem}[Rows at real roots of both signs]\label{thm:bothsigns}
Let $F\in\mathbb C[x]$ be nonzero of degree $D$, and let
$a_1\ge\cdots\ge a_p>0$ and $c_1\ge\cdots\ge c_n>0$, repetitions allowed, with
$p+n\ge1$, be such that $\prod_{i}(x-a_i)\prod_i(x+c_i)$ divides $F$.  Let
$0<\rho<1$ with $\rho a_i>1$ and $\rho c_i>1$ for every $i$.  For
$1\le k\le\max(p,n)$ put
\[
  \Pi_k(\rho)=\prod_{i=k}^{p}(\rho a_i)\prod_{i=k}^{n}(\rho c_i),
\]
an empty product being $1$.  Then
\begin{equation}\label{eq:bothrowsbound}
  b_k(F)\ \ge\ |f_D|\,\frac{\Pi_k(\rho)-1}{(1-\rho)^{-k}-1}.
\end{equation}
\end{theorem}

\begin{proof}
First let $F$ be real.  Since $k\le\max(p,n)\le p+n\le D$, there are at least
$k-1$ positions below $D$.  Let $z_1<\cdots<z_{k-1}$ be the backward distances
$D-j$ of $k-1$ positions $j<D$ carrying the $k-1$ largest nonleading
magnitudes; they are distinct positive integers, so $z_i\ge i$.  Let
$F^*(t)=t^DF(1/t)=\sum_{d=0}^Df_{D-d}t^d$ and
\[
  g=T_{z_{k-1}}\cdots T_{z_1}F^*,\qquad
  [t^d]g=f_{D-d}\prod_{i=1}^{k-1}\Bigl(1-\frac d{z_i}\Bigr).
\]
So $g(0)=f_D\ne0$ and $[t^{z_i}]g=0$.  As in the proof of
Proposition~\ref{prop:rowone}, $(x-\alpha)\mid F$ gives $(1-\alpha t)\mid F^*$,
so $F^*$ has the zeros $1/a_1\le\cdots\le1/a_p$ in $(0,\infty)$ and the zeros
$-1/c_i$, of absolute values $1/c_1\le\cdots\le1/c_n$, in $(-\infty,0)$, with
multiplicity.  Every intermediate polynomial $T_{z_s}\cdots T_{z_1}F^*$ has
constant term $f_D\ne0$, so Lemma~\ref{lem:rolleside} applies at each step;
after $s$ steps the $j$-th positive zero is at most the $(j+1)$-th positive
zero of the previous stage, hence at most $1/a_{j+s}$.  So $g$ has positive
zeros $\eta_j\le1/a_{j+k-1}$, $1\le j\le p-k+1$, and negative zeros
$\eta'_j$ with $|\eta'_j|\le1/c_{j+k-1}$, $1\le j\le n-k+1$, all of modulus
less than $\rho$.  Jensen's formula for $g$ on $|t|\le\rho$ gives
\[
  \log|f_D|+\log\Pi_k(\rho)\ \le\
  \log|f_D|+\sum_{g(\zeta)=0,\ |\zeta|<\rho}\log\frac\rho{|\zeta|}\ \le\
  \log\max_{|t|=\rho}|g(t)|,
\]
because $\log(\rho/|\eta_j|)\ge\log(\rho a_{j+k-1})$, likewise for the
$\eta'_j$, and the remaining terms are positive.  On the other side, for
$d\ge1$ either $[t^d]g=0$ or the position $D-d$ is not among the $k-1$ largest,
so $|f_{D-d}|\le b_k(F)$, while
$\prod_i|1-d/z_i|\le\prod_{i=1}^{k-1}(1+d/i)=\binom{d+k-1}{k-1}$.  Hence
\[
  \max_{|t|=\rho}|g(t)|\ \le\ |f_D|+b_k(F)\sum_{d\ge1}\binom{d+k-1}{k-1}\rho^d
  =|f_D|+b_k(F)\bigl((1-\rho)^{-k}-1\bigr),
\]
and the two displays give \eqref{eq:bothrowsbound}.  For complex $F$ apply the
real case to $\operatorname{Re}(F/f_D)$, as in \cite[Corollary~3.3]{coefficient-mass}: it
is real of degree $D$ with leading coefficient $1$, divisible by the real
polynomial $\prod_i(x-a_i)\prod_i(x+c_i)$, and its coefficients are dominated
by those of $F/f_D$.
\end{proof}

The exclusions cost one root of each sign, and in the worst case the largest
one: $\Pi_k$ omits $a_1,\ldots,a_{k-1}$ and $c_1,\ldots,c_{k-1}$.  This is the
opposite of \cite[Theorem~3.2]{coefficient-mass}, which omits the smallest roots, so
\eqref{eq:bothrowsbound} is the right tool when the moduli are comparable, and
it is sharp there.

\begin{corollary}[Mass at real roots of both signs]\label{cor:bothsignsmass}
Let $F\in\mathbb C[x]$ be nonzero with $|f_D|\ge1$, divisible by
$\prod_{i=1}^L(x-r_i)$ with $r_i$ real, repetitions allowed, and $|r_i|\ge R\ge4$.
Then for $1\le k\le\lceil L/2\rceil$,
\[
  b_k(F)\ \ge\ |f_D|\,\frac{(R/2)^{L-2k+2}-1}{2^k-1},
\]
and
\begin{equation}\label{eq:bothmass}
  \Lambda(F)\ \ge\ \Bigl(\Bigl\lceil\frac L2\Bigr\rceil+1\Bigr)\log|f_D|
  +\Bigl\lfloor\frac{(L+1)^2}4\Bigr\rfloor\log\frac R2
  -\frac12\Bigl\lceil\frac L2\Bigr\rceil\Bigl(\Bigl\lceil\frac L2\Bigr\rceil+3\Bigr)\log2 .
\end{equation}
\end{corollary}

\begin{proof}
Let $p$ and $n$ count the positive and negative $r_i$; then
$\lceil L/2\rceil\le\max(p,n)$.  Take $\rho=\frac12$ in
Theorem~\ref{thm:bothsigns}: every $\rho|r_i|\ge R/2\ge2$, and
$\Pi_k(\tfrac12)\ge(R/2)^{(p-k+1)^++(n-k+1)^+}\ge(R/2)^{L-2k+2}$.  This is the row
bound.  For $k\le\lceil L/2\rceil$ the exponent $L-2k+2$ is at least $1$, so
$X=(R/2)^{L-2k+2}\ge2$, $X-1\ge X/2$, and
$\log b_k(F)\ge\log|f_D|+(L-2k+2)\log(R/2)-(k+1)\log2$.  The leading coefficient
and the $\lceil L/2\rceil$ largest nonleading magnitudes, all nonzero, sit at
distinct positions, so $\Lambda(F)\ge\log^+|f_D|+\sum_k\log^+b_k(F)$ is at
least the sum of these lower bounds; finally
$\sum_{k=1}^{\lceil L/2\rceil}(L-2k+2)=\lfloor(L+1)^2/4\rfloor$ and
$\sum_{k=1}^{\lceil L/2\rceil}(k+1)=\frac12\lceil L/2\rceil(\lceil L/2\rceil+3)$.
\end{proof}

\begin{proposition}[Sharpness of the constant $\frac14$]\label{prop:bothsignssharp}
Let $m\ge1$, $R\ge4$, $\lambda\ge1$ and $s_1,\ldots,s_m\in[R,\lambda R]$.  The
polynomial $P=\prod_{j=1}^m(x^2-s_j^2)$ is a monic multiple of the $L=2m$ real
roots $\pm s_j$, all of modulus at least $R$, and
\[
  \Lambda(P)\ \le\ \Bigl\lfloor\frac{(L+1)^2}4\Bigr\rfloor\log(\lambda R)+m^2\log2 .
\]
So the coefficient $\lfloor(L+1)^2/4\rfloor\sim\frac14L^2$ of $\log R$ in
\eqref{eq:bothmass} cannot be increased, while for roots of one sign
\cite[Corollary~3.4]{coefficient-mass} gives $\frac{L(L+1)}2\log(R-1)$.
\end{proposition}

\begin{proof}
The nonleading coefficients of $P$ are $\pm e_j(s_1^2,\ldots,s_m^2)$,
$1\le j\le m$, each between $1$ and $\binom mj(\lambda R)^{2j}$.  So
$\Lambda(P)\le\sum_j\log\binom mj+m(m+1)\log(\lambda R)$, and
$m(m+1)=\lfloor(2m+1)^2/4\rfloor$, $\sum_j\log\binom mj\le m^2\log2$.
\end{proof}


\section{Thin sectors}\label{sec:thin}

Without separated moduli, a multiple may have a single tropical root
(Theorem~\ref{thm:archnewton}) for all the prescribed roots of an annulus, as $x^{2M}+\rho^{2M}$ and
$x^{2m}+x+B$ of Proposition~\ref{prop:noangle} do.  The angle mechanism of
Proposition~\ref{prop:noangle} puts $K$ roots within an angle $\delta$ of a
ray only at the price of degree about $K/\delta$.  A multiple of smaller
degree is handled by the real theory, applied to the real polynomial
$r\mapsto\operatorname{Im}\bigl(e^{-i\theta}F(re^{i\psi})\bigr)$ along a ray
$\arg z=\psi$ next to the sector.  For $0<x\le\pi/2$ put
\[
  E(x)=\frac{4x}{\log2}\Bigl(1+\log\Bigl(1+\frac{\pi}{2x}\Bigr)\Bigr).
\]
$E$ is increasing, since
$\frac{d}{dx}\bigl(x(1+\log(1+\frac{\pi}{2x}))\bigr)=1+\log(1+t)-\frac t{1+t}>0$
with $t=\pi/(2x)$.

\begin{lemma}[Crossings next to a sector]\label{lem:sectorcount}
Let $B\in\mathbb C[x]$ be nonzero of degree $d$, let $\varphi\in\R$,
$0<\delta'\le\pi/2$ and $R>0$, and suppose that $B$ has no zero other than $0$
on the two rays $\arg z=\varphi\pm\delta'$, and at least $K$ zeros, counted
with multiplicity, in $\{z:|z|\ge R,\ |\arg z-\varphi|<\delta'\}$.  Then there are
$r_1\in[R/2,R)$ and $\psi\in\{\varphi-\delta',\varphi+\delta'\}$ such that for
every $\theta\in\R$ the real polynomial
\[
  h_\theta(r)=\operatorname{Im}\bigl(e^{-i\theta}B(re^{i\psi})\bigr)
  =\sum_{k=0}^d\operatorname{Im}\bigl(e^{-i\theta}b_ke^{ik\psi}\bigr)\,r^k
\]
has at least $K-dE(\delta')/\pi-1$ distinct zeros in $(r_1,\infty)$.
\end{lemma}

\begin{proof}
For $r>0$ let $\Gamma_r=\{re^{it}:|t-\varphi|\le\delta'\}$, traversed in the
direction of increasing $t$, and for a point $\beta$ with $|\beta|\ne r$ let
$\Delta_\beta(r)$ be the change of a continuous argument of $z-\beta$ along
$\Gamma_r$.  We claim
\[
  |\Delta_\beta(r)|\ \le\ \min\Bigl(2\pi,\ \frac{2\delta'r}{\bigl|r-|\beta|\bigr|}\Bigr).
\]
The second bound holds because $|d\arg(z-\beta)|\le|dz|/|z-\beta|$, the arc
has length $2\delta'r$, and $|z-\beta|\ge|r-|\beta||$ on it.  For the first:
if $|\beta|<r$, the argument of $re^{it}-\beta$ has derivative
$(r^2-\operatorname{Re}(\bar\beta re^{it}))/|re^{it}-\beta|^2>0$ and increases by
$2\pi$ over a full turn, so $0<\Delta_\beta(r)<2\pi$; if $|\beta|>r$, the
points $z-\beta$, $|z|=r$, lie in a closed disc not containing $0$, which is
seen from $0$ under an angle less than $\pi$, so $|\Delta_\beta(r)|<\pi$.

Next, for every $\varrho\ge0$ and $a>0$,
\[
  \int_a^{2a}\min\Bigl(2\pi,\frac{2\delta'r}{|r-\varrho|}\Bigr)\frac{dr}r
  \ \le\ 4\delta'\Bigl(1+\log\Bigl(1+\frac\pi{2\delta'}\Bigr)\Bigr)
  =E(\delta')\log2 .
\]
Indeed, the integrand is $\min(2\pi/r,2\delta'/|r-\varrho|)$.  Put
$\eta=\delta'a/\pi$.  Where $|r-\varrho|\le\eta$ it is at most $2\pi/a$ on a
set of measure at most $2\eta$, which contributes at most $4\delta'$.  Elsewhere
it is at most $2\delta'/|r-\varrho|$ on a set of measure at most $a$ on which
$|r-\varrho|>\eta$; such an integral is largest for the set
$\eta<|r-\varrho|\le\eta+a/2$, where it equals
$4\delta'\log(1+a/(2\eta))=4\delta'\log(1+\pi/(2\delta'))$.

Let $\beta_1,\ldots,\beta_d$ be the zeros of $B$ with multiplicity, and
$S(r)=\sum_l\min(2\pi,2\delta'r/|r-|\beta_l||)$.  By the last display,
$\int_a^{2a}S(r)\,dr/r\le dE(\delta')\log2$, while
$\int_a^{2a}dr/r=\log2$; so $S(r)\le dE(\delta')$ on a subset of $[a,2a]$ of
positive measure.  Taking $a=R/2$, choose $r_1\in[R/2,R)$ there, different from
every $|\beta_l|$; taking $a=M$ larger than every $|\beta_l|$, choose $r_2$ in
the same way in $[M,2M]$.  Since
$\arg B(z)=\arg b_d+\sum_l\arg(z-\beta_l)$, the change $\Delta_{\rm arc}(r_i)$
of $\arg B$ along $\Gamma_{r_i}$ satisfies
$|\Delta_{\rm arc}(r_i)|\le S(r_i)\le dE(\delta')$.

Let $\psi_\pm=\varphi\pm\delta'$ and let $\Delta_\pm$ be the change of
$\arg B(re^{i\psi_\pm})$ as $r$ increases from $r_1$ to $r_2$.  The region
$W=\{r_1<|z|<r_2,\ |\arg z-\varphi|<\delta'\}$ has no zero of $B$ on its
boundary and contains the $K$ zeros of the statement (their moduli are at
least $R>r_1$ and below $M<r_2$).  The argument principle along
$\partial W$, traversed counterclockwise, gives
\[
  2\pi K\ \le\ \Delta_-+\Delta_{\rm arc}(r_2)-\Delta_+-\Delta_{\rm arc}(r_1)
  \ \le\ |\Delta_-|+|\Delta_+|+2dE(\delta').
\]
Let $\psi$ be the ray with the larger $|\Delta_\pm|$, so that
$|\Delta_\psi|\ge\pi K-dE(\delta')$, and let $A(r)$ be a continuous argument of
$B(re^{i\psi})$ on $[r_1,r_2]$.  Fix $\theta$.  The open interval between
$A(r_1)$ and $A(r_2)$ has length $|\Delta_\psi|$, so it contains at least
$|\Delta_\psi|/\pi-1$ points of $\theta+\pi\mathbb Z$; by continuity each of
them equals $A(r)$ for some $r\in(r_1,r_2)$, and distinct points give distinct
$r$.  At such $r$ the number $e^{-i\theta}B(re^{i\psi})$ is real, so
$h_\theta(r)=0$.  This gives at least $|\Delta_\psi|/\pi-1\ge K-dE(\delta')/\pi-1$
distinct zeros of $h_\theta$ in $(r_1,r_2)$.
\end{proof}

\begin{theorem}[Thin sectors]\label{thm:thinsector}
Let $B\in\mathbb C[x]$ have degree $d$ and $|b_d|\ge1$, let $\varphi\in\R$,
$0<\delta\le\pi/4$ and $R\ge6$, and suppose that $B$ has at least $K$ zeros,
counted with multiplicity, in $\{z:|z|\ge R,\ |\arg z-\varphi|\le\delta\}$.  If
$d\,E(2\delta)\le\pi K/2$, then with $m=\lceil K/2\rceil-1$,
\[
  b_k(B)\ \ge\ |b_d|\Bigl(\frac R2-1\Bigr)^{m-k+1}\quad(1\le k\le m),
\]
\[
  \Lambda(B)\ \ge\ (m+1)\log|b_d|+\frac{m(m+1)}2\log\Bigl(\frac R2-1\Bigr),
\]
and $\frac{m(m+1)}2\ge\frac{K(K-2)}8$.
\end{theorem}

\begin{proof}
Choose $\delta'\in(\delta,2\delta]$ such that no zero of $B$ lies on the rays
$\arg z=\varphi\pm\delta'$; all but finitely many values qualify.  The $K$
zeros lie in the open sector $|\arg z-\varphi|<\delta'$, and $E(\delta')\le
E(2\delta)$ because $E$ is increasing.  Lemma~\ref{lem:sectorcount} gives
$r_1\ge R/2$ and a ray $\psi$ such that every $h_\theta$ has at least
$K-dE(2\delta)/\pi-1\ge K/2-1$ distinct zeros in $(r_1,\infty)$, hence at least
$m$ of them, say $\varrho_1<\cdots<\varrho_m$, all larger than $R/2\ge3$.
Choose $\theta$ with $e^{-i\theta}b_de^{id\psi}=i|b_d|$.  Then $h_\theta$ is a
nonzero real polynomial of degree $d$ with leading coefficient $|b_d|$ and
coefficients $h_k=\operatorname{Im}(e^{-i\theta}b_ke^{ik\psi})$, so
$|h_k|\le|b_k|$, and it is divisible by $\prod_{i\le m}(x-\varrho_i)$ with
$\varrho_1\ge2$.  \cite[Theorem~3.2]{coefficient-mass} and \cite[Corollary~3.4]{coefficient-mass} give
$b_k(h_\theta)\ge|b_d|\prod_{i=k}^m(\varrho_i-1)$ and
$\Lambda(h_\theta)\ge(m+1)\log|b_d|+\sum_{i=1}^mi\log(\varrho_i-1)$, and
$\varrho_i-1>R/2-1$.  The coefficientwise domination $|h_k|\le|b_k|$, with the
leading coefficients at the same position $d$, gives $b_k(B)\ge b_k(h_\theta)$
and $\Lambda(B)\ge\Lambda(h_\theta)$.  Finally $m+1\ge K/2$ and $m\ge K/2-1$.
\end{proof}

The theorem uses no integrality, and its degree hypothesis cannot be
dropped: $x^{2M}+\rho^{2M}$ in Proposition~\ref{prop:noangle} has about
$2\delta M/\pi$ roots within $\delta$ of the imaginary axis and
$b_2=0$, and its degree $2M$ exceeds $\pi K/(2E(2\delta))$.  At Gaussian roots
the large-degree case is paid for by integrality, as the next section shows.


\section{Gaussian-integer roots}\label{sec:gaussian}

For roots in $\mathbb Z[i]$ and integer multiples there is one more piece of
information: an integer polynomial takes Gaussian-integer values at them,
and a nonzero Gaussian integer has modulus at least one.  This charges the
degree of a multiple against its mass (Theorem~\ref{thm:gausscharge}), which
removes the degree hypothesis of Theorem~\ref{thm:thinsector}.

Throughout, $\alpha_1,\ldots,\alpha_K\in\mathbb Z[i]$ are distinct, with
$\alpha_j=a_j+ic_j$ and $c_j\ge1$; put $\rho_j=|\alpha_j|$,
$\rho_{\min}=\min_j\rho_j$ and
\[
  P=\prod_{j=1}^K(x-\alpha_j)(x-\bar\alpha_j)
   =\prod_{j=1}^K\bigl(x^2-2a_jx+a_j^2+c_j^2\bigr)\in\mathbb Z[x],
\]
a monic polynomial with $2K$ distinct roots and $P(0)\ne0$.  For
$F=\sum_{i=0}^Df_ix^i$ and $0\le m\le D+1$ write
\[
  R_m=\sum_{i<m}f_ix^i,\qquad G_m=\sum_{i\ge m}f_ix^{i-m},\qquad
  F=R_m+x^mG_m .
\]
We call $m\in\{1,\ldots,D\}$ a \emph{split} of $F$.  A split is \emph{clean
at} $\alpha$ if $G_m(\alpha)=0$, and \emph{$P$-clean} if it is clean at every
$\alpha_j$.  A polynomial $B\in\mathbb Z[x]$ is \emph{$P$-irreducible} if
$B(0)\ne0$ and no split of $B$ is $P$-clean.

\begin{lemma}[Gaussian carries]\label{lem:gausscarry}
Let $F\in\mathbb Z[x]$ have degree $D$, let $\alpha\in\mathbb Z[i]$ with
$F(\alpha)=0$, $\rho=|\alpha|$, and let $1\le m\le D$.  Then
$G_m(\alpha)\in\mathbb Z[i]$ and $R_m(\alpha)=-\alpha^mG_m(\alpha)$.  If the
split $m$ is not clean at $\alpha$, then $|R_m(\alpha)|\ge\rho^m$, and if
moreover $\rho>2$, some $i<m$ has
\[
  |f_i|\ \ge\ (\rho/2)^{m-i}.
\]
\end{lemma}

\begin{proof}
$G_m$ has integer coefficients, so $G_m(\alpha)\in\mathbb Z[i]$, and
$0=F(\alpha)=R_m(\alpha)+\alpha^mG_m(\alpha)$.  If $G_m(\alpha)\ne0$ it is a
nonzero Gaussian integer, so $|G_m(\alpha)|\ge1$ and
$|R_m(\alpha)|=\rho^m|G_m(\alpha)|\ge\rho^m$.  If every $i<m$ had
$|f_i|<(\rho/2)^{m-i}$, then
$|R_m(\alpha)|\le\sum_{i<m}|f_i|\rho^i<\rho^m\sum_{i<m}2^{i-m}<\rho^m$, a
contradiction.
\end{proof}

The lemma is the Gaussian form of the gap principle for lacunary
polynomials~\cite{lenstra-lacunary}: across a run of small coefficients a
root of $F$ must be a root of both parts.  Here it is used at every split,
not only at long gaps.

\begin{lemma}[Blocks]\label{lem:gaussblocks}
Every nonzero $F\in\mathbb Z[x]$ divisible by $P$ can be written
\[
  F=\sum_{s=1}^rx^{e_s}B_s,\qquad e_s+\deg B_s<e_{s+1},
\]
with every $B_s\in\mathbb Z[x]$ $P$-irreducible and divisible by $P$.
Consequently $\Lambda(F)=\sum_s\Lambda(B_s)$, $|\operatorname{lead}B_s|\ge1$,
and $b_k(F)\ge b_k(B_s)$ for every $s$ and $k$.
\end{lemma}

\begin{proof}
Write $F=x^eF_0$ with $F_0(0)\ne0$; since $P(0)\ne0$, $P\mid F_0$.  Induct
on $\deg F_0$.  If no split of $F_0$ is $P$-clean, $F_0$ itself is the only
block.  Otherwise let $m$ be a $P$-clean split: $G_m(\alpha_j)=0$ for all
$j$, hence $R_m(\alpha_j)=F_0(\alpha_j)-\alpha_j^mG_m(\alpha_j)=0$, and by
reality of the coefficients also $R_m(\bar\alpha_j)=G_m(\bar\alpha_j)=0$.
Since $P$ is monic in $\mathbb Z[x]$ with $2K$ distinct roots, division with
remainder in $\mathbb Z[x]$ shows $P\mid R_m$ and $P\mid G_m$ in
$\mathbb Z[x]$.  Both are nonzero ($R_m(0)=F_0(0)$, and $G_m$ has the leading
coefficient of $F_0$) and of smaller degree than $F_0$; $R_m(0)\ne0$, and
$G_m$ is handled after removing a power of $x$.  Their supports lie in
$[0,m-1]$ and $[m,\deg F_0]$, so the decompositions given by induction
concatenate.  The supports of the $x^{e_s}B_s$ are disjoint and cover that of
$F$, which gives the additivity of $\Lambda$ and the inclusion of the
nonleading magnitudes of each $B_s$ among those of $F$ (the leading
coefficient of $B_s$, $s<r$, is a nonleading coefficient of $F$).
\end{proof}

\begin{theorem}[Degree charging]\label{thm:gausscharge}
Assume $\rho_{\min}>2$, and let $B\in\mathbb Z[x]$ be $P$-irreducible and
divisible by $P$, of degree $d$.  Then
\[
  \sum_{i<d}\log^+|b_i|\ \ge\ d\log\frac{\rho_{\min}}2,
  \qquad\text{so}\qquad
  \Lambda(B)\ \ge\ \log|b_d|+d\log\frac{\rho_{\min}}2 .
\]
Hence every nonzero $F\in\mathbb Z[x]$ divisible by $P$ has
$\Lambda(F)\ge\sum_s(\deg B_s)\log(\rho_{\min}/2)$ over its blocks.
\end{theorem}

\begin{proof}
Every split $m\in\{1,\ldots,d\}$ of $B$ is unclean at some $\alpha_j$, with
$\rho_j\ge\rho_{\min}$, so by Lemma~\ref{lem:gausscarry} there is $i<m$ with
$|b_i|\ge(\rho_{\min}/2)^{m-i}$.  For a fixed $i$, the splits charged to it
this way have distinct $m-i\ge1$ with
$(m-i)\log(\rho_{\min}/2)\le\log|b_i|$, so there are at most
$\log^+|b_i|/\log(\rho_{\min}/2)$ of them.  All $d$ splits are charged, which
is the first inequality.  The rest follows from $\log^+|b_d|\ge\log|b_d|$ and
Lemma~\ref{lem:gaussblocks}.
\end{proof}

Compare Proposition~\ref{prop:noangle}: $x^{2M}+\rho^{2M}$ pays
$\Lambda=2M\log\rho$ for its degree $2M$, which the theorem allows; what the
theorem forbids, at Gaussian roots, is a block of large degree and small
mass.

\begin{theorem}[Thin Gaussian clusters]\label{thm:thingauss}
Let $\varphi\in\R$, $0<\delta\le\pi/4$ and $R\ge6$, and suppose that every
$\alpha_j$ has $|\alpha_j|\ge R$ and $|\arg\alpha_j-\varphi|\le\delta$.  Every
nonzero $F\in\mathbb Z[x]$ divisible by $P$ satisfies $\Lambda(F)\ge2K\log R$
and
\[
  \Lambda(F)\ \ge\ \min\Bigl\{\frac{K(K-2)}8\log\Bigl(\frac R2-1\Bigr),\
  \frac{\pi K\log(R/2)}{2E(2\delta)}\Bigr\},
  \qquad
  \frac{\pi}{2E(2\delta)}\ \ge\ \frac{\pi\log2}{16\,\delta\log(e\pi/(2\delta))}.
\]
If moreover some block $B$ of $F$ (Lemma~\ref{lem:gaussblocks}) has
$\deg B\cdot E(2\delta)\le\pi K/2$, then
$b_k(F)\ge(R/2-1)^{m-k+1}$ for $1\le k\le m=\lceil K/2\rceil-1$.
\end{theorem}

\begin{proof}
Let $B$ be a block of $F$, of degree $d$.  It is divisible by $P$, so it
vanishes at the $K$ distinct $\alpha_j$; $B(0)\ne0$, $|b_d|\ge1$,
$\Lambda(F)\ge\Lambda(B)$ and $b_k(F)\ge b_k(B)$.  If $dE(2\delta)\le\pi K/2$,
Theorem~\ref{thm:thinsector} applies to $B$ and gives the first term, and the
rows.  Otherwise $d>\pi K/(2E(2\delta))$, and Theorem~\ref{thm:gausscharge},
with $\rho_{\min}\ge R>2$, gives
$\Lambda(B)\ge d\log(\rho_{\min}/2)>\pi K\log(R/2)/(2E(2\delta))$.  Next,
$B=PQ$ with $Q\in\mathbb Z[x]$ since $P$ is monic, so $B(0)=P(0)Q(0)$ is a
nonzero multiple of $P(0)=\prod_j\rho_j^2\ge R^{2K}$, and $B(0)$ is a
coefficient of $F$; this gives $\Lambda(F)\ge2K\log R$.  Finally, for
$0<\delta\le\pi/4$ we have $1+\pi/(4\delta)\le\pi/(2\delta)$, so
$E(2\delta)\le(8\delta/\log2)(1+\log(\pi/(2\delta)))
=(8\delta/\log2)\log(e\pi/(2\delta))$.
\end{proof}


\section{Sectors of every width}\label{sec:widesectors}

The loss $dE(\delta')/\pi$ in Lemma~\ref{lem:sectorcount} is what confines
Theorem~\ref{thm:thinsector} to thin sectors.  This section replaces it by
$\delta'n/\pi$, where $n$ counts only the zeros of modulus above the inner
arc, plus a degree term that becomes small when the inner radius is averaged
over a long range of scales.  The rate $\delta'/\pi$ is the share of a
sector of opening $2\delta'$ in a full turn, and $x^N-\rho^N$ shows that no
smaller rate is possible.  $K$ zeros of modulus at least $R$ in a sector of
half-width $\delta<\pi/2$ then cost either about $\frac12m^2\log R$, where
$m$ is the excess of $K$ over the share $\delta n/\pi$, or about $n\log R$
through the Mahler measure.  So every multiple of degree $O(K)$ pays
$(\pi/\delta-o(1))K\log R$, and this constant is attained
(Theorem~\ref{thm:widesector}, Corollary~\ref{cor:wideconst},
Proposition~\ref{prop:widesharp}); at Gaussian roots the degree restriction
is removed by Theorem~\ref{thm:gausscharge} (Corollary~\ref{cor:widegauss}).
In wide sectors a bound $\Lambda\ge cK^2\log R$ is therefore false for real
multiples, the correct order being $\min(K^2,K/\delta)\log R$.

For $0<\delta\le\pi/2$ put
\[
  \kappa(\delta)=\frac1\pi\int_0^1\min\Bigl(2\arcsin q,\ \frac{2\delta q}{1-q}\Bigr)\frac{dq}q .
\]
The integrand increases with $\delta$ and is at most $2\arcsin(q)/q$, so
$\kappa$ is increasing and, by dominated convergence, continuous.  Since
$\int_0^1\arcsin(q)\,dq/q=\int_0^{\pi/2}s\cot s\,ds=-\int_0^{\pi/2}\log\sin s\,ds
=\frac\pi2\log2$, we have $\kappa(\delta)\le\log2$; and bounding
$2\arcsin q\le\pi q$ and splitting at $q=1-2\delta/\pi$,
\[
  \kappa(\delta)\ \le\ \frac1\pi\Bigl(\int_0^{1-2\delta/\pi}\frac{2\delta\,dq}{1-q}
  +\int_{1-2\delta/\pi}^1\pi\,dq\Bigr)
  =\frac{2\delta}\pi\Bigl(1+\log\frac\pi{2\delta}\Bigr).
\]

\begin{lemma}[Sharp crossings next to a sector]\label{lem:sharpcross}
Let $B\in\mathbb C[x]$ be nonzero of degree $d$, let $\varphi\in\R$ and
$0<\delta'\le\pi/2$, and suppose that $B$ has no zero other than $0$ on the
rays $\arg z=\varphi\pm\delta'$.  Let $0<a<b$ and
$\varepsilon=\kappa(\delta')d/\log(b/a)$.  There is $r_1\in[a,b)$, different
from the modulus of every zero of $B$, with the following property.  Let $n$
be the number of zeros of $B$ of modulus greater than $r_1$, and $N$ the
number of those in the open sector $|\arg z-\varphi|<\delta'$, both counted
with multiplicity.  For $\psi\in\R$ and $\theta\in\R$ put
$h_\theta(r)=\operatorname{Im}\bigl(e^{-i\theta}B(re^{i\psi})\bigr)$.
There is $\psi\in\{\varphi-\delta',\varphi+\delta'\}$ such that for
every $\theta$ the polynomial $h_\theta$ has at least
$N-\delta'n/\pi-\varepsilon-1$ distinct zeros in $(r_1,\infty)$.
\end{lemma}

\begin{proof}
For $r>0$ let $\Gamma_r$ and $\Delta_\beta(r)$ be as in the proof of
Lemma~\ref{lem:sectorcount}, let $[\,\cdot\,]$ be the indicator of a
condition, and for $\beta\ne0$ with $|\beta|\ne r$ put
$q=\min(|\beta|/r,r/|\beta|)<1$.  We claim
\[
  \bigl|\Delta_\beta(r)-2\delta'[\,|\beta|<r\,]\bigr|\ \le\
  c(q)=\min\Bigl(2\arcsin q,\ \frac{2\delta'q}{1-q}\Bigr).
\]
If $|\beta|<r$, write $re^{it}-\beta=re^{it}(1-w)$ with
$w=(\beta/r)e^{-it}$; if $|\beta|>r$, write $re^{it}-\beta=-\beta(1-w)$ with
$w=(r/\beta)e^{it}$.  In both cases $|w|=q<1$, so $1-w$ stays in the disc of
radius $q$ about $1$, where the principal argument is continuous and at most
$\arcsin q$ in absolute value.  So $\Delta_\beta(r)-2\delta'[\,|\beta|<r\,]$
is the change of $\operatorname{Arg}(1-w)$ along $\Gamma_r$, which is at most
$2\arcsin q$ in absolute value; and since $dw/dt=\mp iw$, the derivative
$\frac{d}{dt}\operatorname{Arg}(1-w)=\operatorname{Im}\frac{\pm iw}{1-w}$ is at most $q/(1-q)$ in
absolute value, over a range of length $2\delta'$.  For $\beta=0$,
$\Delta_0(r)=2\delta'$ exactly.

Let $\beta_1,\ldots,\beta_d$ be the zeros of $B$ with multiplicity (all but
the zero at $0$ are in the sum below) and
$S(r)=\sum_{\beta_l\ne0}c(q_l(r))$.  For each $\beta_l\ne0$,
\[
  \int_{-\infty}^{\infty}c\bigl(q_l(e^u)\bigr)\,du
  =2\int_0^\infty c(e^{-v})\,dv=2\int_0^1c(q)\frac{dq}q=2\pi\kappa(\delta'),
\]
so $\int_{\log a}^{\log b}S(e^u)\,du\le2\pi\kappa(\delta')d$ and
$S(e^u)\le2\pi\varepsilon$ on a subset of $[\log a,\log b)$ of positive
measure.  Choose $r_1=e^u$ there, different from every $|\beta_l|$.

Put $X_1=N-\delta'n/\pi-\varepsilon-1$ and choose $\eta>0$
with $\lceil X_1-\eta\rceil=\lceil X_1\rceil$.  Since $c(q)\le\pi q$, $S(r)\to0$ as
$r\to\infty$; choose $r_2$ larger than $r_1$ and every $|\beta_l|$ with
$S(r_2)\le2\pi\eta$.  Let $\psi_\pm=\varphi\pm\delta'$ and let $\Delta_\pm$ be
the change of a continuous argument of $B(re^{i\psi_\pm})$ as $r$ increases
from $r_1$ to $r_2$.  The region
$W=\{r_1<|z|<r_2,\ |\arg z-\varphi|<\delta'\}$ has no zero of $B$ on its
boundary and contains exactly $N$ zeros.  With $\Delta_{\rm arc}(r)$ as in
the proof of Lemma~\ref{lem:sectorcount}, the argument principle gives
$2\pi N=\Delta_-+\Delta_{\rm arc}(r_2)-\Delta_+-\Delta_{\rm arc}(r_1)$.  By
the claim,
$\Delta_{\rm arc}(r)=\sum_l\Delta_{\beta_l}(r)=2\delta'\#\{l:|\beta_l|<r\}+\epsilon(r)$
with $|\epsilon(r)|\le S(r)$, and
$\#\{l:|\beta_l|<r_2\}-\#\{l:|\beta_l|<r_1\}=n$.  Hence
\[
  |\Delta_-|+|\Delta_+|\ \ge\ |\Delta_--\Delta_+|
  \ \ge\ 2\pi N-2\delta'n-S(r_1)-S(r_2)
  \ \ge\ 2\pi N-2\delta'n-2\pi\varepsilon-2\pi\eta .
\]
By the last paragraph of the proof of Lemma~\ref{lem:sectorcount}, for every
$\theta'$ there are at least $|\Delta_\pm|/\pi-1$ distinct
$r\in(r_1,r_2)$ at which $e^{-i\theta'}B(re^{i\psi_\pm})$ is real.

Let $\psi$ be the ray with the larger $|\Delta_\pm|$.  Then
$|\Delta_\psi|/\pi-1\ge X_1-\eta$, so every $h_\theta$ has at least
$\lceil X_1-\eta\rceil=\lceil X_1\rceil\ge X_1$ distinct zeros in
$(r_1,\infty)$.
\end{proof}

\begin{theorem}[Sectors of every width]\label{thm:widesector}
Let $B\in\mathbb C[x]$ have degree $d$ and $|b_d|\ge1$, let $\varphi\in\R$,
$0<\delta<\pi/2$, $T>1$ and $R\ge3T$, and suppose that $B$ has at least $K$
zeros, counted with multiplicity, in
$\{z:|z|\ge R,\ |\arg z-\varphi|\le\delta\}$.  Put
$e=\kappa(\delta)d/\log T$ and, for $\nu\ge0$,
$g(\nu)=\nu\log(R/T)-\log(\nu+1)-1$.  For every integer $m\ge0$,
\[
  \Lambda(B)\ \ge\ \min\Bigl\{\frac{m(m+1)}2\log\Bigl(\frac RT-1\Bigr),\
  g(\nu_m)\Bigr\},\qquad
  \nu_m=\frac\pi\delta\,(K-1-m-e)^+ .
\]
\end{theorem}

\begin{proof}
Choose $\delta'\in(\delta,\pi/2)$ such that no zero of $B$ other than $0$ lies
on the rays $\arg z=\varphi\pm\delta'$; all but finitely many values qualify.
Apply Lemma~\ref{lem:sharpcross} with $a=R/T$ and $b=R$, so
$\varepsilon'=\kappa(\delta')d/\log T$; the $K$ zeros have modulus at least
$R>r_1$ and lie in the open sector, so $N\ge K$.  Note $r_1\ge R/T\ge3$, and
$g$ is increasing on $[0,\infty)$, as $g'(\nu)=\log(R/T)-1/(\nu+1)>0$.

\emph{Mahler measure.}  Suppose $n\ge1$ and let
$B^*(t)=t^dB(1/t)=\sum_kb_{d-k}t^k$, so $B^*(0)=b_d$ and the nonzero zeros of
$B^*$ are the $1/\beta$, $\beta\ne0$ a zero of $B$.  Jensen's formula on
$|t|=\varrho=n/(n+1)\ge\frac12$, where the $n$ points $1/\beta$ with
$|\beta|>r_1$ lie in $|t|<1/3$ and each contributes $\log(\varrho|\beta|)\ge\log(\varrho r_1)$,
and every other contribution is nonnegative, gives
\[
  \log\max_{|t|=\varrho}|B^*(t)|\ \ge\ \log|b_d|+n\log(\varrho r_1)\ \ge\ n\log r_1-1,
\]
using $n\log\frac{n+1}n\le1$.  On the other side
$\max_{|t|=\varrho}|B^*|\le\max_i|b_i|\sum_k\varrho^k=(n+1)\max_i|b_i|$.  So
$\Lambda(B)\ge\log\max_i|b_i|\ge n\log r_1-\log(n+1)-1\ge g(n)$.

\emph{Crossings.}  Choose $\theta$ with $e^{-i\theta}b_de^{id\psi}=i|b_d|$.  As
in the proof of Theorem~\ref{thm:thinsector}, $h=h_\theta$ is real of degree
$d$ with leading coefficient $|b_d|\ge1$ and $|h_k|\le|b_k|$, so
$\Lambda(B)\ge\Lambda(h)$.

Let $X_1=N-\delta'n/\pi-\varepsilon'-1$.  If $X_1\ge m$, then by
Lemma~\ref{lem:sharpcross} $h$ has $m$ distinct zeros larger than
$r_1\ge3$, and \cite[Theorem~3.2]{coefficient-mass} with \cite[Corollary~3.4]{coefficient-mass} gives
$\Lambda(h)\ge\sum_{i=1}^mi\log(\varrho_i-1)\ge\frac{m(m+1)}2\log(R/T-1)$.
Otherwise $n>(\pi/\delta')(K-1-m-\varepsilon')$.  If the right side is positive,
$n\ge1$ and $\Lambda(B)\ge g(n)\ge g(\nu_m(\delta'))$, where $\nu_m(\delta')$
is $\nu_m$ with $\delta'$ and $\varepsilon'$ in place of $\delta$ and $e$; if not,
$\nu_m(\delta')=0$ and $g(0)=-1<0\le\Lambda(B)$.  So the bound holds with
$\nu_m(\delta')$, for all admissible $\delta'$ arbitrarily close to $\delta$;
as $\delta'\downarrow\delta$, $\kappa(\delta')\to\kappa(\delta)$ and
$\nu_m(\delta')\to\nu_m$, and $g$ is continuous.

\end{proof}

The first term is the thin-sector mechanism, now at every width; the second
is the angle mechanism.  If the zeros of $B$ of modulus above $R/T$ are few,
for instance if $B=PQ$ with $P$ the product of $K$ conjugate pairs in the
sector and all zeros of $Q$ in $|z|\le R/T$, then $n\le2K$ and the
theorem with the largest admissible $m$ gives
$\Lambda(B)\ge\frac12\bigl((1-\frac{2\delta}\pi)K-e-2\bigr)^2\log(R/T-1)$.
So in a sector of any
width short of a half-plane a cheap multiple needs about
$(\frac\pi\delta-2)K$ further zeros of modulus comparable to $R$, which the
Mahler term charges linearly.

\begin{corollary}[The constant $\pi/\delta$]\label{cor:wideconst}
In Theorem~\ref{thm:widesector} suppose $d\le\lambda K$ with $\lambda>0$, let
$0<\varepsilon\le\frac12$, and take $T=\exp(\kappa(\delta)\lambda/\varepsilon)$,
assuming $R\ge3T$.  Then
\[
  \Lambda(B)\ \ge\ \min\Bigl\{\frac{\varepsilon^2K^2}2\log\Bigl(\frac RT-1\Bigr),\
  \frac\pi\delta\bigl((1-2\varepsilon)K-2\bigr)\log\frac RT
  -\log\Bigl(\frac{\pi K}\delta+1\Bigr)-1\Bigr\}.
\]
Consequently, for fixed $\delta$ and $\lambda$,
$\Lambda(B)\ge(\pi/\delta-o(1))K\log R$ as $K\to\infty$ and $R\to\infty$.
\end{corollary}

\begin{proof}
Here $e\le\kappa(\delta)\lambda K/\log T=\varepsilon K$.  Take
$m=\lceil\varepsilon K\rceil\le\varepsilon K+1$: then
$\frac{m(m+1)}2\ge\frac{\varepsilon^2K^2}2$ and
$K-1-m-e\ge(1-2\varepsilon)K-2=:\nu_0\delta/\pi$.  Since $g$ is increasing,
$g(\nu_m)\ge g(\nu_0^+)$, and $g(\nu_0)\ge\nu_0\log(R/T)-\log(\pi K/\delta+1)-1$
because $\nu_0\le\pi K/\delta$; if $\nu_0<0$ the claimed bound is negative.
For the limit take $\varepsilon=\max(K^{-1/3},(\log R)^{-1/2})$: then
$\log T=O((\log R)^{1/2})$, the first term is at least
$\frac12K^{4/3}\log(R/T-1)$, which exceeds the second for large $K$, and the second is
$(\pi/\delta)(1-o(1))K\log R$.
\end{proof}

\begin{proposition}[Sharpness of $\pi/\delta$]\label{prop:widesharp}
Let $\varphi\in\R$, $0<\delta\le\pi$, $K\ge1$, $\rho\ge1$ and
$N=\lceil\pi K/\delta\rceil$.  The monic polynomial $x^N-\rho^N$ has at least
$K$ zeros in $\{z:|z|=\rho,\ |\arg z-\varphi|\le\delta\}$, degree
$N\le\pi K/\delta+1$, and
\[
  \Lambda(x^N-\rho^N)=N\log\rho\ \le\ \Bigl(\frac{\pi K}\delta+1\Bigr)\log\rho .
\]
It is real, and integral when $\rho^N\in\mathbb Z$.  If $0<\varphi-\delta$
and $\varphi+\delta<\pi$, it is a real multiple of $K$ conjugate pairs
$\beta_j,\bar\beta_j$ with $\beta_j$ in the sector.
\end{proposition}

\begin{proof}
The zeros $\rho e^{2\pi il/N}$ have arguments spaced by $2\pi/N$, so a closed
arc of length $2\delta$ contains at least $\lfloor\delta N/\pi\rfloor\ge K$ of
them.  The only nonleading nonzero coefficient is $-\rho^N$.  Under the last
hypothesis the $K$ zeros are nonreal, and their conjugates are zeros of the
real polynomial $x^N-\rho^N$.
\end{proof}

Its degree is at most $(\pi/\delta+1)K$, so Corollary~\ref{cor:wideconst}
applies with $\lambda=\pi/\delta+1$: the constant $\pi/\delta$ of the angle
mechanism is sharp in every sector.  The thin-sector term is sharp too: on a
common ray, $\prod_j(x-\rho_je^{i\varphi})$ with $\rho_j\in[R,2R]$ has
$\Lambda\le\frac{K(K+1)}2\log(2R)+K^2\log2$.  So for complex multiples of
degree $O(K)$ the infimum of $\Lambda$ is of order $\min(K^2,K/\delta)\log R$
as $R\to\infty$, with both constants attained in the limit.  In a wide sector, with $\delta$
fixed, it is linear in $K$, and
$\Lambda\ge cK^2\log R$ fails already for the monic real multiple $x^N-\rho^N$.
\begin{corollary}[Gaussian roots in sectors of every width]\label{cor:widegauss}
In the setting of Section~\ref{sec:gaussian}, let $\varphi\in\R$,
$0<\delta<\pi/2$, $T>1$, $R\ge3T$ and $\lambda>0$, suppose that every
$\alpha_j$ has $|\alpha_j|\ge R$ and $|\arg\alpha_j-\varphi|\le\delta$, and put
$e=\kappa(\delta)\lambda K/\log T$.  Every nonzero $F\in\mathbb Z[x]$ divisible
by $P$ satisfies, for every integer $m\ge0$,
\[
  \Lambda(F)\ \ge\ \min\Bigl\{\lambda K\log\frac R2,\
  \frac{m(m+1)}2\log\Bigl(\frac RT-1\Bigr),\ g(\nu_m)\Bigr\},
\]
with $g$ and $\nu_m$ as in Theorem~\ref{thm:widesector}.  In particular, for fixed
$\delta$, $\Lambda(F)\ge(\pi/\delta-o(1))K\log R$ as $K\to\infty$ and
$R\to\infty$.
\end{corollary}

\begin{proof}
Let $B$ be a block of $F$ (Lemma~\ref{lem:gaussblocks}), of degree $d$; it is
real, $|b_d|\ge1$, it vanishes at the $K$ distinct $\alpha_j$, and
$\Lambda(F)\ge\Lambda(B)$.  If $d\ge\lambda K$, Theorem~\ref{thm:gausscharge}
with $\rho_{\min}\ge R>2$ gives $\Lambda(B)\ge\lambda K\log(R/2)$.  Otherwise
$\kappa(\delta)d/\log T\le e$, and Theorem~\ref{thm:widesector} applies to $B$;
its bound only improves when $e$ decreases, since $\nu_m$ then
increases and $g$ is increasing.  For the last statement take
$\lambda=\pi/\delta$ and argue as in Corollary~\ref{cor:wideconst}.
\end{proof}


\paragraph{Data availability.}
No datasets were generated or analyzed in this study.

\begin{thebibliography}{99}
\small
\bibitem{bbbp-ray}
F.~Beaucoup, P.~Borwein, D.~W. Boyd and C.~Pinner,
\newblock Power series with restricted coefficients and a root on a given ray,
\newblock \emph{Mathematics of Computation} 67 (1998), 715--736,
\newblock \href{https://doi.org/10.1090/S0025-5718-98-00960-0}
{doi:10.1090/S0025-5718-98-00960-0}.

\bibitem{coefficient-mass}
B.~Pham,
\newblock Displaced-zero certificates and the coefficient mass of polynomial
multiples,
\newblock manuscript in preparation.

\bibitem{dj}
A.~Dubickas and J.~Jankauskas,
\newblock On the reduced height of a polynomial,
\newblock \emph{Publicationes Mathematicae Debrecen} 71 (2007), 325--348,
\newblock \href{https://doi.org/10.5486/PMD.2007.3728}
{doi:10.5486/PMD.2007.3728}.

\bibitem{lenstra-lacunary}
H.~W. Lenstra, Jr.,
\newblock Finding small degree factors of lacunary polynomials,
\newblock in \emph{Number Theory in Progress}, Vol.~1 (Zakopane, 1997),
de Gruyter, 1999, 267--276.

\bibitem{schinzel-complex}
A.~Schinzel,
\newblock The reduced length of a polynomial with complex coefficients,
\newblock \emph{Acta Arithmetica} 133 (2008), 73--81,
\newblock \href{https://doi.org/10.4064/aa133-1-5}
{doi:10.4064/aa133-1-5}.
\end{thebibliography}

\end{document}
